\chapter{Grommets}

\section*{Definition and Overview}
Grommets, also called ventilation tubes or tympanostomy tubes, are tiny hollow cylinders inserted through the eardrum to ventilate the middle ear and prevent the build-up of fluid. They are most commonly used to treat glue ear (otitis media with effusion) in children, when fluid persists for months and causes hearing loss, and to reduce the frequency of recurrent acute ear infections. Grommets are inserted during a short operation under general anaesthetic, and they usually fall out on their own after 6--18 months, by which time the ear problems have often resolved. They are among the most common surgical procedures performed in children and are safe and effective.

\section*{Epidemiology}
Grommets are one of the most common surgical procedures in childhood. In Australia, around 20,000--30,000 grommet operations are performed each year, most in children under five. The procedure is more common in boys and in children with a history of recurrent ear infections, childcare attendance, and exposure to cigarette smoke. Grommet insertion rates vary considerably between regions, reflecting differences in referral and surgical practice. Most children who have grommets need them only once, though around 20--30\% need a second set.

\section*{Aetiology and Pathophysiology}
Grommets work by addressing the dysfunction that causes middle ear fluid and infections.

\begin{itemize}
    \item \textbf{Eustachian tube dysfunction:} poor drainage of the middle ear leads to fluid accumulation (glue ear) and a breeding ground for infection
    \item \textbf{Ventilation:} the grommet creates an opening in the eardrum, allowing air into the middle ear and equalising pressure
    \item \textbf{Drainage:} fluid can escape through the tube, and the middle ear lining dries out
    \item \textbf{Hearing:} restoring air-filled middle ear space improves sound transmission, correcting the conductive hearing loss of glue ear
    \item \textbf{Infection prevention:} in recurrent acute otitis media, ventilation reduces the conditions that promote infection
    \item \textbf{Spontaneous extrusion:} the grommet migrates outward with the eardrum as it grows, typically falling out within 6--18 months, and the eardrum hole heals
\end{itemize}

\section*{Clinical Features}
Children are offered grommets when they have specific problems.

\begin{itemize}
    \item Glue ear persisting for more than three months with significant hearing loss
    \item Hearing loss affecting speech, language, or learning
    \item Recurrent acute ear infections, typically three or more in six months, or four or more in a year
    \item Ear infections that are severe or complicated, such as those with bulging eardrums or perforation
    \item Symptoms of hearing difficulty: not responding, turning up the television, or speech delay
    \item Balance problems related to middle ear fluid
\end{itemize}

\section*{Differential Diagnosis}
The decision to insert grommets is made after considering other approaches and diagnoses.

\begin{itemize}
    \item Watchful waiting, as many cases of glue ear resolve spontaneously within three months
    \item Hearing aids for persistent glue ear in some cases
    \item Adenoidectomy, which may be performed alongside or instead of grommets
    \item Treatment of underlying allergies or reflux that contribute to ear problems
    \item Speech therapy for language delay that persists after hearing is restored
    \item Investigation for sensorineural hearing loss if conductive causes are excluded
    \item Assessment for other causes of recurrent infections, such as immune problems
\end{itemize}

\section*{When to Seek Medical Care}
Discuss concerns about glue ear, hearing, or recurrent ear infections with your GP, who can arrange hearing tests and referral to an ear, nose, and throat (ENT) specialist.

\textbf{Call 000 (or local emergency services) for:}
\begin{itemize}
    \item Severe ear pain, high fever, or discharge from the ear with a grommet in place
    \item Sudden hearing loss
    \item Signs of meningitis: high fever, stiff neck, severe headache, or sensitivity to light
    \item A grommet that falls out into the ear canal and cannot be removed, or any foreign body in the ear
\end{itemize}

\section*{Diagnosis and Investigations}
The need for grommets is confirmed by examination and hearing tests.

\begin{itemize}
    \item \textbf{Otoscopy:} examination of the eardrum showing fluid behind it or signs of recurrent infection
    \item \textbf{Tympanometry:} confirms the presence of fluid and eardrum dysfunction
    \item \textbf{Hearing assessment:} audiometry quantifies the hearing loss
    \item \textbf{Documentation of infection frequency:} a history of recurrent acute otitis media
    \item \textbf{ENT assessment:} specialist evaluation confirms the indication and plans surgery
    \item \textbf{Pre-operative checks:} general health assessment before anaesthetic
\end{itemize}

\section*{Management}
Grommets are inserted under general anaesthetic as a short day procedure.

\begin{itemize}
    \item \textbf{The operation:} a small incision is made in the eardrum, fluid is suctioned out, and the grommet is placed; the procedure takes 10--15 minutes
    \item \textbf{Recovery:} children go home the same day and are usually back to normal within a day or two
    \item \textbf{Hearing improvement:} hearing usually improves immediately
    \item \textbf{Water precautions:} advice varies, but many surgeons recommend earplugs for swimming and bathing to reduce infection risk
    \item \textbf{Follow-up:} regular reviews, including hearing tests, monitor the grommets and eardrum
    \item \textbf{Spontaneous loss:} grommets fall out on their own within 6--18 months
    \item \textbf{Reinsertion:} a second set is considered if ear problems return
    \item \textbf{Management of discharge:} a discharging grommet is treated with antibiotic drops and review
\end{itemize}

\section*{Complications}
\begin{itemize}
    \item Ear discharge (otorrhoea) through the grommet, which occurs in around 10--20\% of children and is treated with drops
    \item Blockage of the grommet
    \item Premature or delayed extrusion of the grommet
    \item A persistent hole in the eardrum that does not heal after the grommet falls out, in a small proportion of children
    \item Tympanosclerosis (scarring of the eardrum), which rarely affects hearing
    \item The need for repeat surgery if problems recur
    \item Risks of general anaesthetic, which are very low for this procedure
\end{itemize}

\section*{Prognosis and Outcomes}
Grommets are highly effective. Hearing returns to normal immediately in most children, and speech and language development usually catch up. Recurrent ear infections are substantially reduced while grommets are in place. Most children need grommets only once, and by the time they fall out, the Eustachian tube has matured and the underlying problem has resolved. Serious complications are rare, and the overall outlook for children treated with grommets is excellent.

\section*{Prevention and Self-Care}
\begin{itemize}
    \item Discuss with your ENT surgeon the recommended water precautions for your child
    \item Use antibiotic ear drops if the ear starts to discharge, as directed
    \item Keep follow-up appointments, including hearing tests
    \item Do not put anything into the ear, including cotton buds
    \item Watch for signs of ear pain, fever, or discharge and seek review promptly
    \item Protect the ears from smoke, which increases ear problems
    \item Keep immunisations up to date, including the pneumococcal vaccine
    \item Encourage language development with reading and talking, and seek speech therapy early if there is delay
\end{itemize}

\section*{Sources and Further Reading}
The following reputable sources provide further information for healthcare students and patients seeking reliable, current advice about grommets.

\begin{itemize}
    \item Australian Society of Otolaryngology Head and Neck Surgery. \textit{Grommets patient information.} https://www.asohns.org.au/
    \item Healthdirect Australia. \textit{Grommets.} https://www.healthdirect.gov.au/
    \item Raising Children Network (Australian Government). \textit{Ear infections and grommets.} https://raisingchildren.net.au/
    \item Royal Children's Hospital Melbourne. \textit{Grommets information for families.} https://www.rch.org.au/
    \item NHS. \textit{Grommets (ventilation tubes).} https://www.nhs.uk/
    \item American Academy of Otolaryngology. \textit{Ear tubes fact sheet.} https://www.enthealth.org/
\end{itemize}
